\begin{abstract}
Calibration-Free Laser-Induced Breakdown Spectroscopy (CF-LIBS) enables standardless multi-element analysis by deriving plasma composition directly from the Saha--Boltzmann framework, but its computational intensity limits real-time deployment.
We address three challenges in advancing CF-LIBS for real-time and high-throughput plasma diagnostics.
First, we develop the first GPU-accelerated CF-LIBS implementation, built on the JAX framework, which provides automatic differentiation for gradient-based plasma parameter retrieval, vectorization for batch spectral synthesis, and just-in-time compilation to GPU.
Five core algorithms are optimized: (1)~Voigt profile evaluation via the Weideman $N\!=\!36$ rational approximation to the Faddeeva function, achieving $76.4\times$ speedup on an NVIDIA V100S GPU at $10^5$ wavelength grid points; (2)~vectorized Boltzmann fitting using closed-form weighted least squares with pad-and-mask batching, achieving $8.8\times$ speedup at 20 elements; (3)~Anderson-accelerated Saha--Boltzmann iteration, reducing the convergence iteration count by $1.6\times$ on average; (4)~softmax compositional closure with exact Jacobian for gradient-based joint optimization; and (5)~a batch forward model producing $10{,}708$ spectra per second at peak throughput---sufficient for real-time analysis at typical 1--20~Hz laser repetition rates with 100-shot accumulation---with an end-to-end pipeline speedup of $13.6\times$ at batch size 1000.
All GPU kernels maintain full numerical precision: the maximum GPU--CPU relative error is $6.8 \times 10^{-8}$ for Voigt profiles and the batch forward model is bit-identical to sequential computation.
Validation on 74~real LIBS spectra from the Aalto University Spectral Library and 6~ChemCam calibration targets confirms machine-precision GPU--CPU parity across all computational kernels.
Second, we present a systematic benchmark of five element identification pathways---ALIAS peak-matching, full-spectrum NNLS decomposition, a novel two-stage hybrid NNLS+ALIAS identifier, Voigt deconvolution followed by ALIAS, and forward-model concentration thresholding---evaluated on 74~Aalto spectra at resolving powers $\mathrm{RP} = 300$--$1100$.
Per-algorithm standalone evaluation reveals complementary strengths: NNLS achieves the highest recall ($R = 0.922$) but low precision ($P = 0.336$), the hybrid NNLS+ALIAS achieves the best single-method precision ($P = 0.578$, $R = 0.731$), and no individual algorithm dominates across all elements.
Third, GPU-accelerated CMA-ES optimization of a 101-parameter multi-method consensus identification system---spanning method selection weights, per-method configuration, 22~per-element identification confidence thresholds, and an evidence combination layer---achieves $F_{\beta=0.5} = 0.914$ on 74~Aalto spectra in 58~seconds on a single V100S GPU with $2{,}048{,}000$ evaluations, a 121\% improvement over the hand-tuned baseline.
The optimized system exploits the complementary spectral information from all five methods: NNLS detection sensitivity for trace elements via the global Boltzmann emissivity envelope, hybrid identification specificity for major elements, and ALIAS for individual-line confirmation, while learning per-element thresholds that nearly veto spectroscopically intractable elements (Mn threshold $= 0.93$, Na $= 0.87$) and maintain low thresholds for reliably detected elements (Fe $= 0.05$, Si $= 0.12$, Al $= 0.07$).
The GPU basis library underpinning all identification pathways is generated in 3.1~minutes on the V100S using exact partition functions computed from 9{,}626~energy levels, covering 83~elements without reliance on polynomial fits.
\end{abstract}
